\section{Evaluation Setup}
\label{sec:eval-setup}

\subsection{Splits}

All headline numbers are computed on the \textbf{v2 in-distribution test split}: 1{,}008 windows held out from the global pool, never seen during any pipeline training. The split is chronological and family-stratified: each of the 27 scenario families contributes windows to train (4{,}701), val (1{,}011), and test (1{,}008), but from disjoint fault-injection runs. This is the regime where past tickets exist for most families; it is what a mature SRE team would observe after a few months of operation.

For the OOD analysis (Section~\ref{sec:g-series}, G8) we use a complementary leave-one-family-out (LOFO) cross-validation: 27 folds, each holding out an entire family. The result is a per-family novelty F1 that tests whether the system memorizes family-specific patterns or learns generic novelty signals.

Time-ordered visibility is enforced everywhere: a test window at timestamp $t$ may retrieve only memory tickets whose \texttt{available\_as\_memory\_from} timestamp is strictly less than $t$, and may not see tickets from its own dataset run (own-run leakage exclusion).

\subsection{Metrics}

\paragraph{Triage detection.} Following the locked panel:

\begin{itemize}
\item \textbf{PR-AUC strict}: area under the precision-recall curve where positives = \texttt{ticket\_worthy} and negatives = $\{$\texttt{borderline}, \texttt{noise}$\}$.
\item \textbf{PR-AUC inclusive}: same but with borderline counted as positive. A more forgiving variant that captures lift on harder windows.
\item \textbf{ROC-AUC}: standard.
\item \textbf{F1 @ FPR $\leq$ 5\%}: operating-point metric for fixed-budget alert review.
\item \textbf{Expected Calibration Error (ECE)}: for calibration validity.
\end{itemize}

\paragraph{Retrieval.}

\begin{itemize}
\item \textbf{Hit@5} (primary headline). Equal to 1 if any gold ticket appears in the top-5, 0 otherwise; averaged over windows with at least one gold.
\item \textbf{MRR (Mean Reciprocal Rank)}: average of $1 / \mathrm{rank}_{\text{first gold}}$, with zero for windows where no gold appears in the top-K.
\item \textbf{Hit@1}: stricter top-1 hit rate.
\item \textbf{Hit@3}: reported for context.
\end{itemize}

\paragraph{Novelty.} For each pipeline that produces an \texttt{is\_novel} flag:

\begin{itemize}
\item \textbf{Novel precision}: of windows flagged novel, what fraction truly have no gold match in memory.
\item \textbf{Novel recall}: of truly-novel windows, what fraction the pipeline correctly flags.
\item \textbf{Novel F1}.
\end{itemize}

\paragraph{Robustness.}

\begin{itemize}
\item \textbf{Distractor confusion rate}: relative Hit@1 drop when distractor tickets are mixed into memory at ratio $r \in \{0, 10, 25, 50\}\%$.
\item \textbf{Hit@5 invariance}: relative Hit@5 drop at the same ratios.
\item \textbf{LOFO F1}: novelty F1 under leave-one-family-out cross-validation, compared to in-distribution F1.
\end{itemize}

\paragraph{Utility.}

\begin{itemize}
\item \textbf{Expected time-to-diagnose} per resolvable incident, computed under a serial-candidate-review simulation: if the gold appears at rank $k \leq K$, diagnosis time is $30k$ seconds; if no hit, diagnosis falls back to a 30-minute manual baseline. Gives reviewers a single actionable number per pipeline.
\end{itemize}

\subsection{Statistical envelope}

Every reported pairwise comparison uses paired bootstrap with 1{,}000 resamples and a fixed seed (42), preserving window identities across pipelines. 95\% bootstrap confidence intervals accompany every headline number, and pairwise deltas carry their own CIs. Stratified metrics are reported across the cross-cutting axes \texttt{scenario\_family}, \texttt{service\_name}, \texttt{is\_hard\_case}, \texttt{is\_novel}, and \texttt{n\_prior\_family\_tickets} (the depth-stratification axis).

\subsection{Pipelines compared}

The locked comparison panel is:

\begin{table}[h]
\centering
\caption{The locked pipeline panel. \texttt{HGB} = Histogram Gradient Boosting on numeric features; \bienc{} = fine-tuned MiniLM-L6-v2 dense retriever; \hrrf{} = Reciprocal Rank Fusion over BM25 + dense + graph; \kgret{} = rule-based Cypher retrieval over the Neo4j knowledge graph; \logseq{} = small transformer over log sequences; \agent{} = three-stage \texttt{hypothesize/retrieve/verify} chain with Qwen 35B; \tch{} = the final hybrid cascade.}
\label{tab:pipeline-panel}
\small
\begin{tabularx}{\linewidth}{l X r}
\toprule
Pipeline & Role & Inference \\
\midrule
\hgb{} & Triage baseline (no retrieval) & 1 ms/window \\
\bienc{} & Best single retriever (top-1) & 5 ms/window \\
\hrrf{} rule & Best single retriever (top-5) & 10 ms/window \\
\hrrf{} LLM & Precision-heavy retriever using LLM-extracted graph & 10 ms/window \\
\logseq{} & Log-sequence family specialist & 5 ms/window \\
\kgret{} & Graph-structural retriever (Neo4j Cypher) & 10 ms/window \\
\agent{} & Three-stage LLM with novelty flag & 30 s/window \\
\midrule
\tch{} & Tiered cascade composing the above & 16 $\mu$s/window (cached) \\
\bottomrule
\end{tabularx}
\end{table}

The \agent{} is the only pipeline using an LLM at inference time. The cascade uses LLM-extracted artifacts (the Neo4j graph, agent labels) heavily at training/preprocessing time, but the runtime composition is classical (logistic regression + RRF + lookup).
